\documentclass[a4paper,12pt]{article}
\setcounter{secnumdepth}{5}
\setcounter{tocdepth}{3}
\input{/usr/share/LaTeX-ToolKit/template.tex}
\begin{document}
\title{Quadric Surface}
\author{沈威宇}
\date{\temtoday}
\titletocdoc
\sct{Quadric Surface}
\ssc{Cylinder}
A cylinder is a surface that consists of all lines (called rulings) that are parallel to a given line and pass through a given plane curve. If the curve is a circle, it's called circular cylinder. If the curve is a ellipse, it's called elliptic cylinder. If the curve is a parabola, it's called parabolic cylinder. If the curve is a hyperbola, it's called hyperbolic cylinder. If the curve is a polygon, it's called prism.
\ssc{Quadric surface}
A quadric surface is the graph of a second-degree equation in three variables, that is,
\[Ax^2+By^2+Cz^2+Dxy+Eyz+Fxz+Gx+Hy+Iz+J=0.\]
By translation and rotation, it can be brought into either of the two standard forms
\[Ax^2+By^2+Cz^2+J=0\]
or
\[Ax^2+By^2+Iz=0\]
